% !TeX program = xelatex
% 本文件可独立编译。历史型号、完整使用日期和人工确认状态需据真实记录补充。
\documentclass[UTF8,a4paper,12pt,fontset=windows]{ctexart}
\usepackage[left=25mm,right=25mm,top=23mm,bottom=23mm]{geometry}
\usepackage{amsmath,amssymb,booktabs,tabularx,array,enumitem,setspace}
\usepackage[hidelinks,unicode]{hyperref}
\setmainfont{Times New Roman}
\setstretch{1.26}
\setlength{\parindent}{2em}
\setlength{\parskip}{0.22em}
\setlength{\emergencystretch}{2em}
\ctexset{
  section={format=\centering\heiti\zihao{3},beforeskip=1.1em,afterskip=0.7em},
  subsection={format=\heiti\zihao{4},beforeskip=0.65em,afterskip=0.35em}
}
\pagestyle{plain}
\setlist[enumerate]{leftmargin=2em,itemsep=0.4em,topsep=0.35em}
\newcolumntype{Y}{>{\raggedright\arraybackslash}X}
\newcommand{\field}[1]{\textbf{#1：}}
\newcommand{\checkitem}[1]{\par\noindent$\square$\quad\textbf{#1}\par}
\hypersetup{pdftitle={AI工具使用详情},pdfsubject={2026年全国大学生数学建模竞赛A题项目AI使用说明},pdfauthor={},pdfcreator={XeLaTeX}}

\begin{document}
\begin{center}
  {\heiti\zihao{1} AI 工具使用详情}\par
  \vspace{0.35em}
  {\small A题：药材烘干过程的热湿迁移与收缩建模}
\end{center}

本说明依据项目现存的 AI 使用记录、求解程序、修订报告及论文整理，具体说明工具介入的环节、输出的用途和核验依据。涉及人工审核的事项按职责列明，未有确认记录的项目保留待确认状态；代表性提示语采用归纳转述，不作为原始对话逐字记录。

\section{工具信息}
\renewcommand{\arraystretch}{1.25}
\noindent\begin{tabularx}{\textwidth}{@{}p{0.24\textwidth}Y@{}}
\toprule
项目 & 填写内容 \\
\midrule
工具名称 & OpenAI Codex、ChatGPT（项目既有记录所列工具）。\\
版本或型号 & 本次说明整理使用 Codex（GPT-6）；历史使用记录未完整保存模型型号，暂不据当前型号推定此前版本。\\
使用日期 & 现有带日期记录可确认 2026-09-11 至 2026-09-12 期间存在相关活动；完整使用起止日期待参赛队据原始记录补充。本说明整理日期为 2026-09-12。\\
使用形式 & 对话讨论、代码辅助实现与调试、通过本地工具运行计算与检查、图文整理及 LaTeX 排版。\\
主要输入 & 赛题及附件、已有模型与程序、计算输出、报错信息、论文草稿，以及用户提供的共享对话和修改意见。\\
\bottomrule
\end{tabularx}

\section{具体使用目的和环节}
\subsection{题意整理与候选模型讨论}
\field{输入}题目对恒定热物性、变物性、全截面达标和半径收缩四个问题的要求，以及几何尺寸、初始状态和物性公式。\field{输出及用途}AI 辅助整理四问的递进关系，讨论 Fourier--Fick 方程、径向有限体积离散、首达事件和移动边界处理的适用条件；用户提供的共享对话也被用于参考上述总体思路。\field{论文位置}第2节“问题分析”、第5节“模型准备”和第6至9节各问模型。

\subsection{附件整理与边界重构}
\field{输入}环境温度、环境水分及半径时序数据。\field{输出及用途}辅助整理读取、单位换算、采样间隔检查和 PCHIP 插值的程序，并讨论附件覆盖区间以外的边界延拓。现行模型在前4小时采用附件序列，之后以 $50\,{}^\circ\mathrm C$ 和 $0.05\,\mathrm{kg/kg}$ 作为环境延拓假设。\field{论文位置}第5.2节“附件数据检查与边界重构”。这些延拓值属于模型假设，原始附件与计算重构值在使用中予以区分。

\newpage
\subsection{数值程序实现、调试与输出一致性检查}
\field{输入}控制方程、边界条件、物性关系、既有 Python 程序及运行结果。\field{输出及用途}AI 辅助实现和检查控制体通量、三对角系统、温湿场 Picard 迭代、全域最大水分首达定位，以及结果文件导出。针对第二、三问相同物性下的轨迹一致性，现行程序统一采用同一条积分轨迹，第三问按规定时刻抽取。

计算采用161个表面加密节点和1秒终算步长；每隔0.1厘米或60秒的输出要求与内部离散精度分别处理。AI 给出的代码及修改建议须经实际运行和数据对照后才能作为结果依据。\field{论文位置}第5.5至5.6节、第6至8节的求解部分及附录B。

\subsection{收缩模型推导与交叉求解辅助}
\field{输入}半径曲线 $R(t)$、干基含水率的定义、固定域映射方程和已有求解代码。\field{输出及用途}AI 参与材料速度与网格速度的关系讨论、移动控制体库存方程推导、验证程序的独立实现、实际计算及结果整理。现行模型明确采用均匀材料收缩闭合，在 $v_s=w$ 的条件下相对输运项为零。

移动网格实现用于检查固定域映射实现的一致性；两者采用相同闭合假设且可作代数对应，因此高度吻合主要支持实现的一致性。COMSOL 对照和二维轴对称计算另用于检查离散实现与中截面近似。\field{论文位置}第9.1至9.2节及第9.3节相关验证小节。均匀材料收缩仍是一项需要实验检验的物理假设。

\subsection{结果图、流程图与图文说明}
\field{输入}实际求解得到的温湿场、首达时刻、半径曲线、验证误差和已有图形。\field{输出及用途}辅助编写绘图代码，调整坐标、单位、图例、字体和多图布局；整理 DrawIO 求解流程及图文解读。新增时空图、达标区域图和收缩截面图均以已有计算数据为基础。\field{论文位置}第5至9节的几何示意、算法流程和数值结果图。

其中，径向解绘成的轴对称截面图属于计算场重建；阈值等值线表示达标区域边界。二者均按各自的计算含义解释。数值图由绘图程序读取结果生成，几何和流程示意图不作为测量数据。

\subsection{论文表述、LaTeX 排版与本说明整理}
\field{输入}模型推导、结果表、论文草稿及修改意见。\field{输出及用途}辅助整理部分正文、精简重复表述、统一术语与符号、修正公式排印、组织表格和交叉引用，并依据现有项目记录撰写本说明。\field{论文位置}第2、5至10节、AI工具使用声明及附录。图形和文字修改后的数值、条件与限定语均以模型和当前结果文件为核对依据。

\newpage
\section{主要提示方式与使用过程}
\subsection{提示策略与迭代方式}
提示通常围绕一个具体环节展开，给出已有公式或代码、期望输出和约束，再依据报错、数据差异或审阅意见逐项修改。其过程可概括为“给出问题与材料---获取候选解释或实现---运行与比较---修正并保留验证记录”。涉及数据的请求强调读取实际附件与结果；涉及模型的请求要求说明假设、变量含义和适用范围；涉及正文的请求要求保持数值与代码一致。

以下提示语为依据项目任务和交付记录归纳的代表性表述，用于说明交互目的和处理过程，\textbf{并非声称保留了对应的逐字对话}。

\subsection{代表性提示与后续处理}
\begin{enumerate}[label=\textbf{（\arabic*）},leftmargin=2.6em]
  \item \textbf{数据与边界处理。}“根据附件的时间列、单位和采样间隔整理边界函数；说明插值区间及区间外采用的假设，检查半径是否产生非物理回弹。”

  对应输出包括数据检查与插值实现。处理时对照原始采样点和时序范围，环境波动予以保留，半径按其单调性质重构；长期延拓假设在正文中单独说明。

  \item \textbf{离散格式与首达事件。}“检查径向有限体积实现的轴线、表面对流和非线性系数更新，区分计算网格与输出位置；所有位置达标时应怎样定位首次时刻？”

  对应输出包括离散检查建议和事件定位代码。现行实现先检验全域最大水分，再确认中心控制；在相邻时间层跨过 $0.15\,\mathrm{kg/kg}$ 时插值定位，并同步生成达标剖面。

  \item \textbf{收缩方程与独立实现。}“从干基水分和移动控制体收支出发，检查材料速度、网格速度及密度定义是否一致，并给出与固定域程序配对的验证实现。”

  对应输出包括条件推导、移动网格程序、比较数据及正文增补。后续保留均匀材料收缩的明确前提，检查均匀场保持、几何守恒和方程平衡，并限制结论适用范围。

  \item \textbf{图文整理与一致性复查。}“用当前终算数据补充时空场和误差图，统一坐标与图例；检查论文表格、正文数值和源数据是否一致，保留必要的物理限定。”

  对应输出包括绘图程序、图表与文字修改。核对时以当前场数据和汇总文件为准，检查图表引用、表内数值和工作簿回读结果；排版通过编译与页面渲染复查。
\end{enumerate}

\subsection{输出取舍原则}
能够说明问题且与附件、程序结果一致的建议，作为实现或表述材料保留；改变方程、边界或解释范围的建议，须连同推导和数值对照讨论。通用方法名称、语言流畅程度和软件间吻合本身，均不能代替模型适用性论证。

\newpage
\section{采纳、人工修改与核验}
下列采纳、修改及程序核验均有项目记录支持。人工执行者及确认时间尚未逐项记录，故按建模、编程和论文整理环节列明复核职责，确认状态见第5节。

\subsection{离散精度与共享轨迹：采纳后修订}
\field{采纳及修改}保留有限体积与隐式迭代框架，明确区分计算网格与输出位置，终算采用161节点、1秒步长，统一第二、三问轨迹。\field{核验依据}空间采用41、81、161、321节点对照，时间采用2秒与1秒对照；第二、三问共享结果完全一致。

\field{结论影响}早期离散结果已由终算结果替代。当前第三、四问达标时间分别为 $57.4719\,\mathrm h$、$51.0877\,\mathrm h$。第三问空间加密和时间加密的相对变化分别约为 $1.65\times10^{-5}$、$1.52\times10^{-5}$，用于评估时空离散收敛性。\field{人工复核职责}编程环节核对设置与输出，建模环节判断精度对结论的影响。

\subsection{收缩机理与效应解释：附条件采纳}
\field{采纳及修改}采纳固定域映射与移动控制体对照的技术路线，明确干基含水率、干骨架连续性及 $v_s=w$ 的闭合前提；将缺少条件的“移动边界项”解释改为有明确速度假设的方程描述。

\field{核验依据}保留均匀材料场、方程收支、移动网格配对及 COMSOL 对照记录；另以问题四相同物性和环境进行固定半径对照。\field{结论影响}归因方式得到修正：问题三和问题四物性不同，不能将两问时长之差全部归因于收缩。验证支持既定假设下的数值一致性，收缩闭合的实验有效性仍待确认。\field{人工复核职责}建模环节复核守恒推导、假设依据与结论范围。

\subsection{正文数值与图表：依据结果文件修正}
\field{采纳及修改}按当前数据修正图文展示值。2026-09-12 的核验记录显示，第三问6小时中心水分由1.0156修正为1.0170，12小时由0.4548修正为0.4562，达标表面水分由0.0525修正为0.0527（单位均为 $\mathrm{kg/kg}$）。

\field{核验依据}该次记录对14张数值表、四份结果工作簿和场数据进行了核对，报告表格错误列表为空、工作簿与源值差异为零。\field{结论影响}上述修改更新了正文展示值，未在该次核验中重算仿真场，亦未改变首达判据。\field{人工复核职责}论文整理环节逐项核对展示值、单位、图注与引用，编程环节复核来源。

\subsection{参赛队的判断与责任}
中截面近似、长期边界、收缩闭合、首达判据和同物性对照均需实质判断。AI 参与按本说明披露；参赛队对方案选择、模型解释和提交内容负责。人工理解与审核情况据实确认，不以程序检查通过代替。

\newpage
\section{人工核验清单}
以下清单供提交前由参赛队逐项确认。现有资料包含程序运行和图文检查记录，尚不足以确认每一项人工复核均已完成，故复选框预留为空。每项确认后可勾选，并在队内保留审核环节、日期及修订依据；本文件不记录姓名、学校或赛区信息。

\checkitem{1）核对公式推导、单位、边界条件和符号一致性。}
核对 $r$、$\xi$、$R(t)$ 的坐标含义，轴线零通量与表面 Robin 条件的方向，摄氏温度与开尔文温度的使用，以及 $C$ 的干基定义。重点复算第四问材料速度、网格速度和干质量连续性之间的关系，区分题给有效热方程与包含潜热的总焓模型。

\textbf{现有依据：}模型准备和各问推导、均匀场检查、方程平衡与移动网格验证记录。

\checkitem{2）在独立环境中运行全部代码并复现论文图表。}
从原始附件开始运行求解、收敛验证、灵敏度、二维与移动网格验证、结果导出及绘图程序，检查第三、四问首达时间，第二、三问共同轨迹，工作簿回读与论文表格。记录运行环境及差异处理结果。

\textbf{现有依据：}项目已有多种离散设置、程序重算和输出比对记录；这些记录本身尚不能证明参赛队已在另一套独立环境中完成全部复现。

\checkitem{3）核验引用、数据来源和事实陈述。}
逐项对照赛题附件与物性公式来源，核实参考文献的题名、作者、年份及所支持的具体内容；区分附件观测、插值重构、模型假设、数值模拟和实验结果。共享 AI 对话只作为使用记录中的参考材料，涉及方法依据的文献须回到原始来源核验。

\textbf{现有依据：}附件检查记录、论文参考文献及各验证报告；文献逐项人工阅读的完成情况待确认。

\checkitem{4）确认核心建模与分析由参赛队主导。}
由参赛队确认实际分工与决策过程，能够独立说明降维依据、边界延拓、非线性物性、阈值事件、收缩闭合和对照方案，解释为何采纳或修改相应建议。原创性说明应对应具体的问题适配、分析与判断，并与已披露的 AI 参与范围一致。

\textbf{确认依据：}队内实际工作和审核记录；仅从程序和论文成品不能判定每一步的人工贡献。

\checkitem{5）确认论文声明与本详情文件一致。}
核对正文“思路讨论、程序实现及文字排版”的声明与本说明一致，并覆盖部分推导、交叉验证、图文增补和本说明整理；补全历史型号及完整日期。检查最终支撑材料是否包含本说明，确认文件中不含参赛者身份信息。

\vspace{0.25em}
\textbf{记录说明：}本说明主要依据项目既有 AI 使用记录、模型修订与验收报告、移动网格交付记录，以及当前数值验证和论文表格核验文件编写。相关文件的相对路径列于同目录说明文件，便于回查。本次工作核对已有记录并编译本说明，未重新运行整套仿真。

\end{document}
